\chapter{2008年安徽卷（文）}

\section{单选题}

\begin{problem}
若 \(A\) 为全体实数的集合，\(B = \{-2, -1, 1, 2\}\)，则下列结论正确的是（\quad）

\choices
    {\(A \cap B = \{-2, -1\}\)}
    {\((\complement_{\R}A) \cup B = (-\infty, 0)\)}
    {\(A \cup B = (0, +\infty)\)}
    {\((\complement_{\R}A)\cap B = \{-2, -1\}\)}
\end{problem}

\begin{answer}
D
\end{answer}

\begin{solution}
按原 PDF 的字面题面，\(A=\R\)，所以 \(A\cap B=B\)，\(\complement_{\R}A=\varnothing\). 因此选项 A、B、C、D 分别变为 \(B\)、\(B\)、\(\R\)、\(\varnothing\)，均不等于题面给出的右侧集合，故字面题面下没有正确选项. 公开同题版本将 \(A\) 印为全体正实数；若按该疑似漏字的版本取 \(A=(0,+\infty)\)，则 \((\complement_{\R}A)\cap B=\{-2,-1\}\)，答案为 D. 现稿保留原 PDF 题面，并记录题面与答案冲突.
\end{solution}

\begin{problem}
若 \(\overrightarrow{AB} = (2, 4)\)，\(\overrightarrow{AC} = (1, 3)\)，则 \(\overrightarrow{BC} =\)（\quad）

\choices
    {\((1,1)\)}
    {\((-1, - 1)\)}
    {\((3,7)\)}
    {\((-3, - 7)\)}
\end{problem}

\begin{answer}
B
\end{answer}

\begin{solution}
由点的位移关系 \(\overrightarrow{BC}=\overrightarrow{AC}-\overrightarrow{AB}\)，得 \(\overrightarrow{BC}=(1,3)-(2,4)=(-1,-1)\)，所以正确选项为 B.
\end{solution}

\begin{problem}
已知 \(m\)，\(n\) 是两条不同直线，\(\alpha\)，\(\beta\)，\(\gamma\) 是三个不同平面，下列命题中正确的是（\quad）

\choices
    {若 \(\alpha \perp \gamma, \beta \perp \gamma\)，则 \(\alpha \parallel \beta\)}
    {若 \(m \perp \alpha\)，\(n \perp \alpha\)，则 \(m \parallel n\)}
    {若 \(m \parallel \alpha\)，\(n \parallel \alpha\)，则 \(m \parallel n\)}
    {若 \(m \parallel \alpha\)，\(m \parallel \beta\)，则 \(\alpha \parallel \beta\)}
\end{problem}

\begin{answer}
B
\end{answer}

\begin{solution}
两条不同直线都垂直于同一平面时，它们的方向都与该平面的法向量平行，因而两直线互相平行，所以选项 B 正确. 选项 A 不成立，例如取 \(\gamma\) 为水平面，取两个相交的竖直平面 \(\alpha\)、\(\beta\)，它们都垂直于 \(\gamma\). 选项 C 不成立，例如取同一水平面上方的两条互相垂直的水平直线，它们都平行于该水平面但彼此不平行. 选项 D 不成立，例如取两个相交的竖直平面，并取一条不在任一平面内且方向平行于它们交线的直线，该直线同时平行于两个平面. 因此只有 B 正确.
\end{solution}

\begin{problem}
\(a < 0\) 是方程 \(ax^2 + 1 = 0\) 有一个负数根的（\quad）

\choices
    {必要不充分条件}
    {充分必要条件}
    {充分不必要条件}
    {既不充分也不必要条件}
\end{problem}

\begin{answer}
B
\end{answer}

\begin{solution}
若 \(a<0\)，则 \(x^2=-1/a>0\)，方程的两个根为 \(\pm\sqrt{-1/a}\)，其中必有一个负根，所以 \(a<0\) 是充分条件. 反之，若方程有负根 \(r<0\)，则 \(ar^2+1=0\)，从而 \(a=-1/r^2<0\)，所以它也是必要条件. 因此该条件既充分又必要，正确选项为 B.
\end{solution}

\begin{problem}
在三角形 \(ABC\) 中，\(AB = 5\)，\(AC = 3\)，\(BC = 7\)，则 \(\angle BAC\) 的大小为（\quad）

\choices
    {\(\frac{2\pi}{3}\)}
    {\(\frac{5\pi}{6}\)}
    {\(\frac{3\pi}{4}\)}
    {\(\frac{\pi}{3}\)}
\end{problem}

\begin{answer}
A
\end{answer}

\begin{solution}
由余弦定理，\(BC^2=AB^2+AC^2-2AB\cdot AC\cos\angle BAC\)，即 \(49=25+9-30\cos\angle BAC\)，所以 \(\cos\angle BAC=-1/2\). 由于三角形内角属于 \((0,\pi)\)，故 \(\angle BAC=2\pi/3\)，正确选项为 A.
\end{solution}

\begin{problem}
函数 \(f(x) = (x - 1)^2 + 1\)（\(x \leqslant 0\)）的反函数为（\quad）

\choices
    {\(f^{-1}(x) = 1 - \sqrt{x - 1}\)（\(x \geqslant 1\)）}
    {\(f^{-1}(x) = 1 + \sqrt{x - 1}\)（\(x \geqslant 1\)）}
    {\(f^{-1}(x) = 1 - \sqrt{x - 1}\)（\(x \geqslant 2\)）}
    {\(f^{-1}(x) = 1 - \sqrt{x - 1}\)（\(x \geqslant 2\)）}
\end{problem}

\begin{answer}
C
\end{answer}

\begin{solution}
令 \(y=f(x)\)，则 \((x-1)^2=y-1\)，所以 \(x=1\pm\sqrt{y-1}\). 由于原函数的定义域为 \(x\leqslant0\)，只能取负号；并且 \(f(x)\geqslant2\)，故反函数的定义域为 \(x\geqslant2\). 因此 \(f^{-1}(x)=1-\sqrt{x-1}\)（\(x\geqslant2\)），对应选项 C. 原卷中 C、D 两个选项的文字相同，故 D 也呈现为同一表达式，但按现稿答案记为 C.
\end{solution}

\begin{problem}
设 \((1 + x)^8 = a_0 + a_1x + \dots + a_8x^8\)，则 \(a_0\)，\(a_1\)，\(\dots\)，\(a_8\) 中奇数的个数为（\quad）

\choices
    {\(2\)}
    {\(3\)}
    {\(4\)}
    {\(5\)}
\end{problem}

\begin{answer}
A
\end{answer}

\begin{solution}
由二项式定理，\(a_k=\mathrm C_8^k\)（\(k=0\)，\(1\)，\(\ldots\)，\(8\)），所以系数依次为 \(1\)，\(8\)，\(28\)，\(56\)，\(70\)，\(56\)，\(28\)，\(8\)，\(1\). 其中只有首尾两个系数为奇数，故奇数的个数为 \(2\)，正确选项为 A.
\end{solution}

\begin{problem}
函数 \( y = \sin \left(2x + \frac{\pi}{3}\right) \) 图象的对称轴方程可能是（\quad）

\choices
    {\(x = -\frac{\pi}{6}\)}
    {\(x = -\frac{\pi}{12}\)}
    {\(x = \frac{\pi}{6}\)}
    {\(x = \frac{\pi}{12}\)}
\end{problem}

\begin{answer}
D
\end{answer}

\begin{solution}
正弦函数在自变量为 \(\frac\pi2+k\pi\)（\(k\in\Z\)）时取得极值，这些竖线是其图象的对称轴. 因此 \(2x+\frac\pi3=\frac\pi2+k\pi\)，即 \(x=\frac\pi{12}+\frac{k\pi}{2}\). 取 \(k=0\) 得 \(x=\frac\pi{12}\)，故正确选项为 D.
\end{solution}

\begin{problem}
设函数 \(f(x) = 2x + \frac{1}{x} - 1\)（\(x < 0\)），则 \(f(x)\)（\quad）

\choices
    {有最大值}
    {有最小值}
    {是增函数}
    {是减函数}
\end{problem}

\begin{answer}
A
\end{answer}

\begin{solution}
在 \(x<0\) 上，\(f'(x)=2-\frac1{x^2}\). 当 \(x<-1/\sqrt2\) 时，\(f'(x)>0\)；当 \(-1/\sqrt2<x<0\) 时，\(f'(x)<0\). 所以函数先增后减，在 \(x=-1/\sqrt2\) 处取得极大值 \(-1-2\sqrt2\). 又 \(f(x)\to-\infty\)（当 \(x\to-\infty\) 或 \(x\to0^-\)），故该极大值就是最大值，正确选项为 A.
\end{solution}

\begin{problem}
若过点 \(A(4,0)\) 的直线 \(l\) 与曲线 \((x - 2)^2 + y^2 = 1\) 有公共点，则直线 \(l\) 的斜率的取值范围为（\quad）

\choices
    {\((- \sqrt{3}, \sqrt{3})\)}
    {\([- \sqrt{3}, \sqrt{3}]\)}
    {\(\left(-\frac{\sqrt{3}}{3}, \frac{\sqrt{3}}{3}\right)\)}
    {\(\left[-\frac{\sqrt{3}}{3}, \frac{\sqrt{3}}{3}\right]\)}
\end{problem}

\begin{answer}
D
\end{answer}

\begin{solution}
圆心为 \(O(2,0)\)，半径为 \(1\). 过 \(A(4,0)\) 的非竖直直线可写为 \(y=k(x-4)\)，即 \(kx-y-4k=0\). 它与圆有公共点，当且仅当圆心到直线的距离不超过半径，即 \(\frac{2|k|}{\sqrt{k^2+1}}\leqslant1\). 平方并整理得 \(3k^2\leqslant1\)，所以 \(k\in[-\sqrt3/3,\sqrt3/3]\). 竖直直线 \(x=4\) 到圆心距离为 \(2>1\)，不产生额外情形，故正确选项为 D.
\end{solution}

\begin{problem}
若 \(A\) 为不等式组 \(\begin{cases}
x \leqslant 0,\\
y \geqslant 0,\\
y - x \leqslant 2
\end{cases}\) 表示的平面区域，则当 \(a\) 从 \(-2\) 连续变化到 \(1\) 时，动直线 \(x + y = a\) 扫过 \(A\) 中的那部分区域的面积为（\quad）

\choices
    {\(\frac{3}{4}\)}
    {\(1\)}
    {\(\frac{7}{4}\)}
    {\(2\)}
\end{problem}

\begin{answer}
C
\end{answer}

\begin{solution}
令 \(s=x+y\)，则 \(y=s-x\)，且变换 \((x,y)\mapsto(s,x)\) 的雅可比绝对值为 \(1\). 三个不等式化为 \(x\leqslant0\)、\(x\leqslant s\)、\(x\geqslant s/2-1\).

当 \(-2\leqslant s\leqslant0\) 时，\(s/2-1\leqslant x\leqslant s\)，截线长度为 \(s/2+1\)；当 \(0\leqslant s\leqslant1\) 时，\(s/2-1\leqslant x\leqslant0\)，截线长度为 \(1-s/2\). 因此
\[
\int_{-2}^{0}\left(\frac{s}{2}+1\right)\,\mathrm ds+\int_{0}^{1}\left(1-\frac{s}{2}\right)\,\mathrm ds=1+\frac34=\frac74
\]
故正确选项为 C.
\end{solution}

\begin{problem}
\(12\) 名同学合影，站成前排 \(4\) 人后排 \(8\) 人，现摄影师要从后排 \(8\) 人中抽 \(2\) 人调整到前排，若其他人的相对顺序不变，则不同调整方法的总数是（\quad）

\choices
    {\(\mathrm{C}_8^2\mathrm{A}_6^6\)}
    {\(\mathrm{C}_8^2\mathrm{A}_3^2\)}
    {\(\mathrm{C}_8^2\mathrm{A}_6^2\)}
    {\(\mathrm{C}_8^2\mathrm{A}_5^2\)}
\end{problem}

\begin{answer}
C
\end{answer}

\begin{solution}
先从后排 \(8\) 人中选出调到前排的 \(2\) 人，有 \(\mathrm C_8^2\) 种. 选定这两人后，原前排 \(4\) 人的相对顺序以及后排剩余 \(6\) 人的相对顺序都保持不变. 把两名调入前排的同学按顺序插入由原前排形成的 \(6\) 个位置，等价于从 \(6\) 个位置中选出两个并安排这两人，有 \(\mathrm A_6^2\) 种. 因此总数为 \(\mathrm C_8^2\mathrm A_6^2\)，对应选项 C.
\end{solution}

\section{填空题}

\begin{problem}
函数 \( f(x) = \frac{\sqrt{|x - 2| - 1}}{\log_2(x - 1)} \) 的定义域为\fillinblank{}.
\end{problem}

\begin{answer}
\([3,+\infty)\)
\end{answer}

\begin{solution}
根式有意义要求 \(|x-2|-1\geqslant0\)，即 \(x\leqslant1\) 或 \(x\geqslant3\). 对数有意义且分母不为零要求 \(x-1>0\) 且 \(\log_2(x-1)\ne0\)，即 \(x>1\) 且 \(x\ne2\). 取交集得定义域为 \([3,+\infty)\).
\end{solution}

\begin{problem}
已知双曲线 \(\frac{x^2}{n} - \frac{y^2}{12 - n} = 1\) 的离心率是 \(\sqrt{3}\). 则 \(n = \)\fillinblank{}.
\end{problem}

\begin{answer}
\(4\)
\end{answer}

\begin{solution}
双曲线的半实轴平方为 \(a^2=n\)，半虚轴平方为 \(b^2=12-n\)，所以焦半距平方为 \(c^2=a^2+b^2=12\). 由离心率 \(c/a=\sqrt3\)，得
\(\frac{12}{n}=3\)，\(n=4\)
且 \(0<n<12\)，该值满足双曲线条件.
\end{solution}

\begin{problem}
在数列 \(\{a_{n}\}\) 中，\(a_{n}=4n-\frac{5}{2}\)，\(a_{1}+a_{2}+\cdots+a_{n}=an^{2}+bn\)，\(n\in\N^{*}\)，其中 \(a\)，\(b\) 为常数，则 \(ab=\)\fillinblank{}.
\end{problem}

\begin{answer}
\(-1\)
\end{answer}

\begin{solution}
由等差数列求和公式，\(a_1+\cdots+a_n=\sum_{k=1}^n\left(4k-\frac52\right)=4\cdot\frac{n(n+1)}2-\frac52n=2n^2-\frac12n\). 比较 \(an^2+bn\) 的系数，得 \(a=2\)、\(b=-1/2\)，从而 \(ab=-1\).
\end{solution}

\begin{problem}
已知 \(A\)，\(B\)，\(C\)，\(D\) 在同一个球面上，\(AB \perp\) 平面 \(BCD\)，\(BC \perp CD\)，若 \(AB = 6\)，\(AC = 2\sqrt{13}\)，\(AD = 8\)，则 \(B\)，\(C\) 两点间的球面距离是\fillinblank{}.
\end{problem}

\begin{answer}
\(\frac{4\pi}{3}\)
\end{answer}

\begin{solution}
设球心为 \(O\)，\(H\) 为 \(O\) 在平面 \(BCD\) 上的射影. 平面 \(BCD\) 截球面所得圆的圆心是 \(H\)，所以 \(HB=HC=HD\). 由于 \(AB\) 垂直于平面 \(BCD\)，设 \(OH=h\)，则 \(A\)、\(B\) 到平面 \(BCD\) 的有向距离分别为 \(6\) 和 \(0\)，从 \(OA=OB\) 得 \(HB^2+(h-6)^2=HB^2+h^2\)，故 \(h=3\). 在直角三角形 \(ABC\)、\(ABD\) 中，\(BC^2=AC^2-AB^2=52-36=16\)，\(BD^2=AD^2-AB^2=64-36=28\). 又 \(BC\perp CD\)，所以三角形 \(BCD\) 在 \(C\) 处直角，\(CD^2=BD^2-BC^2=12\). 其外接圆半径为斜边的一半，即 \(r=BD/2=\sqrt7\). 若球半径为 \(R\)，由截面圆半径公式 \(r^2=R^2-3^2\)，得 \(R=4\).

设 \(B\)、\(C\) 所对的劣弧对应的球心角为 \(\theta\). 弦长公式给出 \(BC=2R\sin\frac{\theta}{2}\)，即 \(4=8\sin\frac{\theta}{2}\). 因为取劣弧，\(0<\theta<\pi\)，所以 \(\theta=\pi/3\). 球面距离是该劣弧的弧长，故为 \(R\theta=4\cdot\frac\pi3=\frac{4\pi}{3}\).
\end{solution}

\section{解答题}

\begin{problem}
已知函数 \( f(x) = \cos \left(2x - \frac{\pi}{3}\right) + 2\sin \left(x - \frac{\pi}{4}\right)\sin \left(x + \frac{\pi}{4}\right) \).

\begin{enumerate}
\item 求函数 \( f(x) \) 的最小正周期；
\item 求函数 \( f(x) \) 在区间 \( \left[-\frac{\pi}{12}, \frac{\pi}{2}\right] \) 上的值域.
\end{enumerate}
\end{problem}

\begin{solution}
\begin{enumerate}
\item 由积化和差公式，\(2\sin\left(x-\frac\pi4\right)\sin\left(x+\frac\pi4\right)=\cos\left(-\frac\pi2\right)-\cos2x=-\cos2x\). 因此 \(f(x)=\cos\left(2x-\frac\pi3\right)-\cos2x=\sin\left(2x-\frac\pi6\right)\). 因为自变量的系数为 \(2\)，最小正周期满足 \(2T=2\pi\)，所以 \(T=\pi\).
\item 当 \(x\in\left[-\frac\pi{12},\frac\pi2\right]\) 时，令 \(t=2x-\frac\pi6\)，则 \(t\in\left[-\frac\pi3,\frac{5\pi}6\right]\)，且 \(f(x)=\sin t\). 该区间包含极大值点 \(\frac\pi2\)，所以最大值为 \(1\). 在该区间内，正弦函数的最小值出现在左端点，等于 \(\sin\left(-\frac\pi3\right)=-\frac{\sqrt3}{2}\)；右端点函数值为 \(\frac12\)，不更小. 因此值域为 \(\left[-\frac{\sqrt3}{2},1\right]\).
\end{enumerate}
\end{solution}

\begin{problem}
在某次普通话测试中，为测试汉字发音水平，设置了 \(10\) 张卡片，每张卡片印有一个汉字的拼音，其中恰有 \(3\) 张卡片上的拼音带有后鼻音“\(\mathrm{g}\)”.

\begin{enumerate}
\item 现对三位被测试者先后进行测试，第一位被测试者从这 \(10\) 张卡片中随机抽取 \(1\) 张，测试后放回，余下 \(2\) 位的测试，也按同样的方法进行，求这三位被测试者抽取的卡片上，拼音都带有后鼻音“\(\mathrm{g}\)”的概率；
\item 若某位被测试者从 \(10\) 张卡片中一次随机抽取 \(3\) 张，求这三张卡片上，拼音带有后鼻音“\(\mathrm{g}\)”的卡片不少于 \(2\) 张的概率.
\end{enumerate}
\end{problem}

\begin{solution}
\begin{enumerate}
\item 每次抽取后放回，所以三次抽取相互独立. 每次抽到带后鼻音“g”的卡片的概率为 \(\frac3{10}\)，故三位被测试者抽到的卡片都带“g”的概率为 \(\left(\frac3{10}\right)^3=\frac{27}{1000}\).
\item 三张卡片中带“g”的卡片数不少于 \(2\) 张，分为恰有 \(2\) 张和恰有 \(3\) 张两种互斥情形. 三张卡片等可能的取法共有 \(\mathrm C_{10}^3\) 种；恰有 \(2\) 张带“g”有 \(\mathrm C_3^2\mathrm C_7^1\) 种，恰有 \(3\) 张带“g”有 \(\mathrm C_3^3\) 种. 因此所求概率为 \(\frac{\mathrm C_3^2\mathrm C_7^1+\mathrm C_3^3}{\mathrm C_{10}^3}=\frac{21+1}{120}=\frac{11}{60}\).
\end{enumerate}
\end{solution}

\begin{problem}
如图，在四棱锥 \(O - ABCD\) 中，底面 \(ABCD\) 四边长为 \(1\) 的菱形，\(\angle ABC = \frac{\pi}{4}\). \(OA \perp\) 底面 \(ABCD\)，\(OA = 2\)，\(M\) 为 \(OA\) 的中点.

\begin{enumerate}
\item 求异面直线 \(AB\) 与 \(MD\) 所成角的大小；
\item 求点 \(B\) 到平面 \(OCD\) 的距离.
\end{enumerate}

\bitmapfigure[max width=0.34\linewidth]{img/2008/anhui_liberal/q19_fig1.png}
\end{problem}

\begin{solution}
\begin{enumerate}
\item 建立直角坐标系：取 \(A=(0,0,0)\)、\(B=(1,0,0)\)，令 \(s=1/\sqrt2\)，则可取 \(D=(-s,s,0)\)、\(C=(1-s,s,0)\)、\(O=(0,0,2)\)、\(M=(0,0,1)\). 于是 \(\overrightarrow{AB}=(1,0,0)\)，\(\overrightarrow{MD}=(-s,s,-1)\). 两向量的长度分别为 \(1\)、\(\sqrt2\)，点积为 \(-s\). 异面直线所成的角取锐角，故其余弦为 \(\frac{|-s|}{1\cdot\sqrt2}=\frac12\)，从而角度为 \(\frac\pi3\).
\item 因为 \(AB\parallel CD\)，且 \(AB\) 不在平面 \(OCD\) 内，所以点 \(A\)、\(B\) 到平面 \(OCD\) 的距离相等. 由 \(\overrightarrow{OC}=(1-s,s,-2)\)、\(\overrightarrow{OD}=(-s,s,-2)\)，可取平面 \(OCD\) 的法向量为 \((0,2,s)\)，其方程为 \(2y+s(z-2)=0\). 点 \(A=(0,0,0)\) 到该平面的距离为 \(\frac{2s}{\sqrt{4+s^2}}=\frac{\sqrt2}{\sqrt{9/2}}=\frac23\)，所以点 \(B\) 到平面 \(OCD\) 的距离为 \(\frac23\).
\end{enumerate}
\end{solution}

\begin{problem}
设函数 \( f(x) = \frac{a}{3} x^3 - \frac{3}{2} x^2 + (a + 1)x + 1 \)，其中 \( a \) 为实数.

\begin{enumerate}
\item 已知函数 \(f(x)\) 在 \(x = 1\) 处取得极值，求 \(a\) 的值；
\item 已知不等式 \( f'(x) > x^2 - x - a + 1 \) 对任意 \( a \in (0, +\infty) \) 都成立，求实数 \( x \) 的取值范围.
\end{enumerate}
\end{problem}

\begin{solution}
\begin{enumerate}
\item \(f'(x)=ax^2-3x+a+1\). 函数在 \(x=1\) 处取得极值，必须有 \(f'(1)=0\)，即 \(2a-2=0\)，所以 \(a=1\). 此时 \(f'(x)=(x-1)(x-2)\)，在 \(x=1\) 左侧为正、右侧为负，确实在 \(x=1\) 处取得极大值.
\item 将不等式移到一边，得 \(f'(x)-\left(x^2-x-a+1\right)=(a-1)x^2-2x+2a=a(x^2+2)-(x^2+2x)>0\). 由于 \(x^2+2>0\)，若 \(x^2+2x\leqslant0\)，则上式等于正数 \(a(x^2+2)\) 与非负数 \(-x^2-2x\) 之和，对任意 \(a>0\) 都大于零. 反之，若 \(x^2+2x>0\)，取充分接近 \(0\) 的正数 \(a\)，上式就会小于零. 因此必要且充分的条件为 \(x^2+2x\leqslant0\)，即 \(-2\leqslant x\leqslant0\).
\end{enumerate}
\end{solution}

\begin{problem}
设数列 \(\{a_{n}\}\) 满足 \(a_1 = a\)，\(a_{n+1} = c a_n + 1 - c\)，\(n \in \N^*\)，其中 \(a\)，\(c\) 为实数，且 \(c \neq 0\).

\begin{enumerate}
\item 求数列 \(\{a_{n}\}\) 的通项公式；
\item 设 \(a = \frac{1}{2}\)，\(c = \frac{1}{2}\)，\(b_n = n(1 - a_n)\)，\(n \in \N^*\)，求数列 \(\{b_n\}\) 的前 \(n\) 项和 \(S_n\)；
\item 若 \(0 < a_{n} < 1\) 对任意 \(n \in \N^{*}\) 成立，证明 \(0 < c \leqslant 1\).
\end{enumerate}
\end{problem}

\begin{solution}
\begin{enumerate}
\item 令 \(u_n=a_n-1\)，则 \(u_{n+1}=ca_n+1-c-1=cu_n\)，所以 \(u_n=(a-1)c^{n-1}\). 因此 \(a_n=1+(a-1)c^{n-1}\).
\item 此时 \(a_n=1-2^{-n}\)，所以 \(b_n=n2^{-n}\). 设 \(S_n=\sum_{k=1}^n k/2^k\). 将 \(S_n/2\) 的求和指标右移一位，作错位相减得
\[
\begin{aligned}
S_n-\frac12S_n
&=\frac12+\sum_{k=2}^{n}\frac{k}{2^k}-\left(\sum_{k=2}^{n}\frac{k-1}{2^k}+\frac{n}{2^{n+1}}\right)\\
&=\frac12+\sum_{k=2}^{n}\frac1{2^k}-\frac{n}{2^{n+1}}=1-\frac{n+2}{2^{n+1}}.
\end{aligned}
\]
由于左边等于 \(S_n/2\)，故 \(S_n=2-\frac{n+2}{2^n}\).
\item 由 \(0<a_1=a<1\) 及 \(a_2=1+c(a-1)=1-c(1-a)<1\)，可知 \(c>0\)，否则 \(a_2\geqslant1\). 若 \(c>1\)，则存在正整数 \(n\) 使 \((1-a)c^{n-1}\geqslant1\)，从而 \(a_n=1-(1-a)c^{n-1}\leqslant0\)，与题设矛盾. 因此 \(0<c\leqslant1\).
\end{enumerate}
\end{solution}

\begin{problem}
设椭圆 \(C: \frac{x^2}{a^2} + \frac{y^2}{b^2} = 1\)（\(a > b > 0\)），其相应于焦点 \(F(2,0)\) 的准线方程为 \(x = 4\).

\begin{enumerate}
\item 求椭圆 \(C\) 的方程；
\item 已知过点 \(F_{1}(-2,0)\) 倾斜角为 \(\theta\) 的直线交椭圆 \(C\) 于 \(A\)，\(B\) 两点，求证：\(|AB| = \frac{4\sqrt{2}}{2 - \cos^2\theta}\)；
\item 过点 \(F_{1}(-2,0)\) 作两条互相垂直的直线分别交椭圆 \(C\) 于 \(A\)，\(B\) 和 \(D\)，\(E\)，求 \(|AB| + |DE|\) 的最小值.
\end{enumerate}
\end{problem}

\begin{solution}
\begin{enumerate}
\item 焦点为 \((2,0)\)，故焦半距 \(c=2\). 准线为 \(x=a^2/c=4\)，所以 \(a^2=8\). 又 \(b^2=a^2-c^2=8-4=4\)，椭圆方程为
\[
\frac{x^2}{8}+\frac{y^2}{4}=1
\]
\item 设直线的单位方向向量为 \((\cos\theta,\sin\theta)\)，其参数方程为 \(x=-2+t\cos\theta\)、\(y=t\sin\theta\). 代入椭圆方程并整理得
\[
(2-\cos^2\theta)t^2-4\cos\theta\,t-4=0
\]
两交点对应的参数根之差的绝对值为
\[
\frac{\sqrt{16\cos^2\theta+16(2-\cos^2\theta)}}{2-\cos^2\theta}
=\frac{4\sqrt2}{2-\cos^2\theta}
\]
由于参数方向是单位向量，根之差的绝对值就是弦长，故
\[
|AB|=\frac{4\sqrt2}{2-\cos^2\theta}
\]
\item 令 \(u=\cos^2\theta\in[0,1]\). 与第一条直线垂直的直线方向角为 \(\theta+\pi/2\)，故其弦长为 \(4\sqrt2/(1+u)\). 两弦长之和为
\[
4\sqrt2\left(\frac1{2-u}+\frac1{1+u}\right)
=\frac{12\sqrt2}{2+u-u^2}
\]
而
\[
2+u-u^2=\frac94-\left(u-\frac12\right)^2\leqslant\frac94
\]
所以和的最小值为
\[
\frac{12\sqrt2}{9/4}=\frac{16\sqrt2}{3}
\]
在 \(u=1/2\) 时取得.
\end{enumerate}
\end{solution}
